\chapter{Conservation of Mechanical Energy}

\section*{Introduction}

A roller coaster is the perfect physical metaphor. At the top of a hill, the car has maximum potential energy and minimum kinetic energy. As it descends, potential energy converts to kinetic energy---the car speeds up. At the bottom of a loop, it is moving fastest (maximum kinetic, minimum potential). If there were no friction, the car would return to exactly the same height at the end of each loop. With friction, each successive hill must be lower---some energy has been ``lost'' to heat. The equation $\frac{1}{2}mv^2 + mgh = \text{const}$ tracks this energy budget precisely.
\section*{Fundamental Equation Used}

For a particle of mass $m$ moving in a conservative force field:
\[
\frac{1}{2}mv^2 + mgh = \text{constant}
\]
or equivalently, $\frac{1}{2}mv_i^2 + mgh_i = \frac{1}{2}mv_f^2 + mgh_f$, relating any two points along the trajectory.
\section*{Assumptions}

\textbf{Assumptions:} Only conservative forces (gravity, spring, electrostatic) do work. Non-conservative forces do no net work or are absent. The system is closed.


\textbf{Limitations:} Does not apply when non-conservative forces are significant (then $\frac{1}{2}mv^2 + mgh + W_{\text{friction}} = \text{const}$). Also does not apply in relativistic mechanics (where kinetic energy is $\gamma mc^2 - mc^2$) or quantum mechanics.


\section*{Mathematical Derivation}

\textbf{Step 1: Start from Newton's second law.}

For a particle of mass $m$ subject to a conservative force $\mathbf{F}$, Newton's second law gives:
\begin{equation}
\mathbf{F} = m\mathbf{a} = m\frac{d\mathbf{v}}{dt}
\end{equation}

\textbf{Step 2: Define a conservative force.}

A force is conservative if it can be written as the gradient of a scalar potential:
\begin{equation}
\mathbf{F} = -\nabla U(\mathbf{r})
\end{equation}
where $U(\mathbf{r})$ is the potential energy. Equivalently, $\nabla \times \mathbf{F} = 0$ (the force field is irrotational).

\textbf{Step 3: Take the dot product with velocity.}

Dot both sides of Newton's law with the velocity $\mathbf{v} = d\mathbf{r}/dt$:
\begin{equation}
\mathbf{F} \cdot \mathbf{v} = m\frac{d\mathbf{v}}{dt} \cdot \mathbf{v}
\end{equation}

\textbf{Step 4: Rewrite both sides as total time derivatives.}

The left side:
\begin{equation}
\mathbf{F} \cdot \mathbf{v} = -\nabla U \cdot \frac{d\mathbf{r}}{dt} = -\frac{dU}{dt}
\end{equation}
by the chain rule. The right side:
\begin{equation}
m\frac{d\mathbf{v}}{dt} \cdot \mathbf{v} = \frac{d}{dt}\left(\frac{1}{2}m\,\mathbf{v} \cdot \mathbf{v}\right) = \frac{dK}{dt}
\end{equation}
where $K = \frac{1}{2}mv^2$ is the kinetic energy.

\textbf{Step 5: Combine and integrate.}

Substituting back:
\begin{equation}
\frac{dK}{dt} = -\frac{dU}{dt}
\end{equation}
Rearranging:
\begin{equation}
\frac{d}{dt}(K + U) = 0
\end{equation}
Integrating over time:
\begin{equation}
K + U = \text{constant}
\end{equation}

\textbf{Step 6: Write the conservation law explicitly.}

For a particle under gravity ($U = mgh$), evaluated at any two points along the trajectory:
\begin{equation}
\boxed{\frac{1}{2}mv_i^2 + mgh_i = \frac{1}{2}mv_f^2 + mgh_f}
\end{equation}
This states that the total mechanical energy is the same at every point along the motion. Note: this result depends critically on the force being conservative. If a non-conservative force (friction) does work $W_{\text{nc}}$, the equation generalizes to $\Delta K + \Delta U = W_{\text{nc}}$.

\section*{Characteristics of the Equation}

\begin{itemize}
\item \textbf{State function:} Mechanical energy depends only on the current state (position and velocity), not on the path.
\item \textbf{Conservative forces:} A force is conservative if $\nabla \times \mathbf{F} = 0$ (irrotational), with $U(\mathbf{r}) = -\int \mathbf{F} \cdot d\mathbf{r}$.
\item \textbf{Work--energy theorem:} $W_{\text{net}} = \Delta K$. For conservative forces only, $W_{\text{cons}} = -\Delta U$, giving $\Delta K + \Delta U = 0$.
\item \textbf{Turning points:} At $v = 0$, all energy is potential. Turning points define the classically allowed region.
\item \textbf{Effective potential:} For central force problems, angular momentum adds a ``centrifugal'' term creating an effective potential for radial motion.
\end{itemize}
\section*{Solution Methods}

\begin{enumerate}
\item \textbf{Energy conservation:} Set $E_i = E_f$ between any two points. Solve for unknowns directly---no need to integrate the equation of motion.
\item \textbf{Energy diagrams:} Plot $U(x)$ and draw a horizontal line at $E$. Intersections are turning points; the gap is kinetic energy.
\item \textbf{Lagrangian mechanics:} Write $L = T - V$ and apply Euler--Lagrange. If $L$ is independent of time, energy is conserved.
\item \textbf{Phase space:} In $(x, p)$, energy conservation defines closed curves revealing the nature of motion.
\end{enumerate}
\section*{Verifications}

\begin{itemize}
\item \textbf{Free fall:} Dropping masses from known heights confirms $v = \sqrt{2gh}$.
\item \textbf{Pendulum timing:} Period measured at various amplitudes matches the energy-conservation prediction, including large-angle corrections.
\item \textbf{Roller coasters:} Accelerometers confirm $K + U \approx \text{const}$, with deficits matching measured friction losses.
\item \textbf{Particle accelerators:} Energy conservation in relativistic collisions is verified to extraordinary precision---missing energy signals new particles.
\end{itemize}
\section*{Applications}

\begin{itemize}
\item \textbf{Pendulum:} Converts between kinetic and potential energy; period found from energy conservation at large amplitudes.
\item \textbf{Orbital mechanics:} Kepler's laws and satellite orbits derived from energy and angular momentum conservation.
\item \textbf{Ballistic motion:} Trajectory of a projectile (neglecting air resistance) is fully determined by energy conservation.
\item \textbf{Spring--mass systems:} Relates amplitude, frequency, and maximum speed of an oscillator.
\item \textbf{Particle physics:} Energy conservation (including rest mass) constrains which reactions are kinematically allowed.
\end{itemize}
\section*{What Richard Feynman Would Have Asked}

\subsection*{1. Plain-English Translation}
\textit{``What does conservation of mechanical energy mean if you had to explain it to a ten-year-old?''}

It means you cannot get something for nothing. If you want a ball to go faster, you have to drop it from higher up. If you want a roller coaster to go through a loop, you have to start from a high enough hill. The total ``score''---kinetic energy (speed) plus potential energy (height)---always adds up to the same number, as long as nothing is rubbing or scraping.

Think of it like money in a bank account with two compartments: a ``cash'' compartment (kinetic energy---how fast you are moving) and a ``savings'' compartment (potential energy---how high up you are). You can move money between compartments, but the total never changes. Friction is like a hidden tax---it takes a small cut from every transaction, and eventually your account runs down.

\subsection*{2. Localized Mechanics}
\textit{``At a single point in space, what does energy conservation look like?''}

At any instant, the particle has a certain kinetic energy ($\frac{1}{2}mv^2$) and potential energy ($mgh$). As the particle moves, these trade back and forth---like a pendulum swinging. At the extremes (highest point), all energy is potential. At the lowest point, all is kinetic. The conservation equation is a constraint: the particle cannot reach a point where the required kinetic energy (from $K + U = E$) is negative---that region is classically forbidden.

He would press you on the microscopic origin. At the molecular level, every collision conserves kinetic energy. The macroscopic conservation is an average over many collisions. When friction converts kinetic energy to heat, it transfers energy from ordered center-of-mass motion to disordered molecular motion.

\subsection*{3. Testing Extreme Limits}
\textit{``Where does energy conservation break down?''}

At the quantum level, energy conservation holds exactly for each interaction, but the Heisenberg uncertainty principle allows temporary violations ($\Delta E \cdot \Delta t \gtrsim \hbar/2$)---virtual particles borrow energy for very short times. These are not violations but quantum phenomena.

In general relativity, energy conservation becomes subtle. In an expanding universe, photons lose energy (cosmological redshift)---the ``missing'' energy goes into the gravitational field of expanding space. For isolated systems in flat spacetime, energy conservation is exact. At the Planck scale, quantum gravity effects may modify the concept, but no experiment has probed this regime.

\subsection*{4. Dimensionality and Geometry}
\textit{``How does the shape of the potential affect energy conservation?''}

The shape of $U(x)$ determines the motion's character. A parabolic potential ($U \propto x^2$) gives simple harmonic motion. A cubic potential gives anharmonic oscillations with amplitude-dependent periods. A double-well potential allows tunneling between minima (in quantum mechanics). In central force problems, the $1/r$ gravitational potential gives closed elliptical orbits; deviations give precessing orbits (general relativity).

\subsection*{5. Cross-Disciplinary Connection}
\textit{``Where else does energy conservation appear outside of physics?''}

Energy conservation is a special case of Noether's theorem: every continuous symmetry gives a conserved quantity. Time-translation symmetry gives energy; space-translation gives momentum; rotation gives angular momentum. The same structure appears in fluid dynamics (conservation of mass, momentum, energy), electrical circuits (conservation of charge and energy), and economics (conservation of money in closed transactions).

The deep reason energy is conserved is that the laws of physics do not change with time. If they were different today than yesterday, energy would not be conserved. The fact that they are the same is one of the deepest observations about the universe.


%=============================================================
% Chapter 18: Continuity Equation
%=============================================================
\section*{FAQ}

\textbf{Q1: If energy is conserved, why do objects eventually stop?}

A: Energy is conserved overall, but mechanical energy (kinetic + potential) is not when friction is present. Kinetic energy converts to thermal energy. The total energy of the universe is constant, but useful mechanical energy decreases.

\textbf{Q2: Can energy conservation ever be violated?}

A: No verified violation exists. In certain exotic situations (cosmological expansion), the definition of ``total energy'' becomes ambiguous, but in all laboratory contexts, energy conservation holds to extraordinary precision.

\textbf{Q3: Why does potential energy depend on the reference level?}

A: Only differences are physically meaningful. You can choose any reference level as long as you are consistent. The choice adds a constant that cancels when computing energy differences.
\section*{Important Bibliography}

\begin{enumerate}
\item Goldstein, H., Poole, C., and Safko, J., \textit{Classical Mechanics}, 3rd ed. (Addison-Wesley, 2002), Chapter 2.
\item Taylor, J.R., \textit{Classical Mechanics}, 2nd ed. (University Science Books, 2005), Chapter 7.
\item Feynman, R.P., Leighton, R.B., and Sands, M., \textit{The Feynman Lectures on Physics}, Vol. I (Addison-Wesley, 1963), Chapter 4.
\item Landau, L.D. and Lifshitz, E.M., \textit{Mechanics}, 3rd ed. (Butterworth-Heinemann, 1976), Chapter 1.
\end{enumerate}


